% math4cs-hw8-trees.tex

% !TEX program = xelatex
%%%%%%%%%%%%%%%%%%%%
% see http://mirrors.concertpass.com/tex-archive/macros/latex/contrib/tufte-latex/sample-handout.pdf
% for how to use tufte-handout
\documentclass[a4paper, justified]{tufte-handout}

\input{hw-preamble} % feel free to modify this file if you understand LaTeX well
%%%%%%%%%%%%%%%%%%%%
\title{计算机数学 (9-trees: 树)}
\me{魏恒峰}{hfwei@hnu.edu.cn}{}{}
\date{2026年5月22日}
%%%%%%%%%%%%%%%%%%%%
\begin{document}
\maketitle
%%%%%%%%%%%%%%%%%%%%
\noplagiarism % PLEASE DON'T DELETE THIS LINE!
%%%%%%%%%%%%%%%%%%%%
\begin{abstract}
\end{abstract}
%%%%%%%%%%%%%%%%%%%%
\beginrequired
%%%%%%%%%%%%%%%

%%%%%%%%%%%%%%%
\begin{problem}[树的顶点度数]
  已知一棵树 $T$ 有 1 个 2 度顶点,
  2 个 3 度顶点, 1 个 4 度顶点, 1 个 5 度顶点,
  其余顶点均为叶节点。

  \noindent 请计算 $T$ 的顶点数。~\footnote{
    提示：上次作业第一题。
  }
\end{problem}

\begin{proof}
\end{proof}
%%%%%%%%%%%%%%%

%%%%%%%%%%%%%%%
\begin{problem}[树的刻画]
  给定无向图 $G$。
  请证明: $G$ 是树当且仅当 $G$ 没有 loop 且 $G$ 有唯一的生成树。
\end{problem}

\begin{proof}
\end{proof}
%%%%%%%%%%%%%%%

%%%%%%%%%%%%%%%
\begin{problem}[生成树与桥]
  给定无向连通图 $G$ 与 $G$ 中的某条边 $e$。
  请证明: $e$ 是桥 (bridge~\footnote{bridge 也称为 cut-edge (割边)。})
    当且仅当 $e$ 属于 $G$ 的每个生成树。
\end{problem}

\begin{proof}
\end{proof}
%%%%%%%%%%%%%%%

%%%%%%%%%%%%%%%
\begin{problem}[最小生成树算法]
  请分别使用 Kruskal 算法与 Prim 算法 (从顶点1开始)
  给出下图的最小生成树。
  要求给出边添加的顺序。

  \fig{width = 0.50\textwidth}{figs/st-example}
\end{problem}

\begin{proof}
\end{proof}
%%%%%%%%%%%%%%%

%%%%%%%%%%%%%%%
\begin{problem}[唯一最小生成树]
  设 $G$ 是无向连通带权图, $T$ 是 $G$ 的一个最小生成树。

  \noindent 请证明: $T$ 是 $G$ 的唯一最小生成树当且仅当
  对于不在 $T$ 中的每一条边 $e$,
  $e$ 的权重大于 $T + e$ 所产生的圈中其它每条边的权重。
\end{problem}

\begin{proof}
\end{proof}
%%%%%%%%%%%%%%%

%%%%%%%%%%%%%%%
\begin{problem}[树的中心点]
  如果 $G$ 是一个连通图，则 $G$ 的一个中心点
  是具有如下性质的顶点 $v$：从 $v$ 到其他所有顶点的距离的最大值尽可能小。

  \noindent 请设计算法, 给定一棵树 $T$, 找出 $T$ 的中心点
  (注意：可能不唯一)，并说明其正确性。
\end{problem}

\begin{proof}
\end{proof}
%%%%%%%%%%%%%%%

\beginoptional
%%%%%%%%%%%%%%%
\begin{problem}[谍影重重]
  \fig{width = 0.70\textwidth}{figs/Bourne}
\end{problem}
%%%%%%%%%%%%%%%

%%%%%%%%%%%%%%%%%%%%
% 如果没有需要订正的题目，可以把这部分删掉
\begincorrection
%%%%%%%%%%%%%%%%%%%%

%%%%%%%%%%%%%%%%%%%%
% 如果没有反馈，可以把这部分删掉
\beginfb

你可以写 (也可以发邮件)
\begin{itemize}
  \item 对课程及教师的建议与意见
  \item 教材中不理解的内容
  \item 希望深入了解的内容
  \item $\cdots$
\end{itemize}
%%%%%%%%%%%%%%%%%%%%
\end{document}